\chapter{引言}
\label{chap:introduction}

\section{研究背景与意义}
\label{sec:background}

本节系统介绍大语言模型、强化学习后训练及其在实际应用中面临的核心挑战，为后续研究工作奠定基础。

\subsection{大语言模型的崛起与能力演进}
\label{subsec:llm_evolution}

近十年间，人工智能研究在深度学习推动下持续演进。其中，大语言模型（Large Language Models，LLMs）取得了快速发展，逐步成为当代AI应用的重要基础设施。从早期的BERT\cite{devlin2018bert}、GPT-2到近年来的GPT-4\cite{openai2023gpt4}、Llama系列\cite{touvron2023llama}、Qwen系列和DeepSeek，这一演进过程中模型能力和规模的增长轨迹极为显著：模型参数规模由亿级（$10^8$）扩展至千亿级（$10^{11}$），训练语料由数百GB扩展至数十TB，对应的计算成本也呈指数级增长。随着规模提升，模型在语言理解、知识表征与文本生成等方面的能力也获得了显著提升。

这场技术演进的起点，是Vaswani等人于2017年提出的Transformer架构\cite{vaswani2017attention}。Transformer以自注意力机制（Self-Attention）取代了循环神经网络（RNN）对序列信息的逐步传递，使得模型在捕捉长程依赖关系上具有了结构性的先天优势。相比RNN需要顺序处理序列，Transformer的并行计算能力使其能够在更短时间内处理更长的上下文，这为大规模预训练的出现奠定了基础。

在此基础上，以\textit{"预测下一个Token"为核心任务的自回归语言模型}（Autoregressive Language Model）确立了大规模预训练的基本范式：通过在海量互联网文本上最小化负对数似然损失，模型被迫学习并压缩了关于人类语言、世界知识乃至逻辑结构的大量信息。这一范式的有效性已被广泛验证：在足够大的数据和计算预算下，模型的语言理解和生成能力往往与训练数据规模呈现幂律关系（Scaling Law）。

\textbf{预训练的局限性与后训练的必要性。} 然而，仅仅通过预训练阶段获得的模型，在实际应用中存在明显的局限性。预训练模型的优化目标是对自然语言分布建模，而非针对特定任务的准确完成。这导致原始的预训练模型在执行用户指令时往往缺乏针对性，甚至会生成偏离预期、内容有害的输出。为了解决这一问题，监督微调（Supervised Fine-Tuning，SFT）技术被引入到LLM的训练流程之中。

监督微调的核心思路是：在人工标注的高质量\textit{"问题—答案"对}数据上，以最大似然估计的方式对模型进行进一步训练，使其学会\textit{"模仿"专家给出的示范性回答}。该方法能够提升模型的指令遵循能力，并使输出更符合人类阅读预期。InstructGPT\cite{ouyang2022training}和相关工作已经证明，适当的SFT能够显著改进模型的实用性。

但是，SFT方法的性能仍然受训练数据质量与覆盖范围的显著约束。对于数学推导、逻辑推理、代码生成等高难度任务，构建规模充足且质量稳定的专家标注数据代价极高。更重要的是，SFT的优化目标仍然主要对应于示范分布拟合，而非显式的推理过程优化。因此，模型往往更擅长复现训练数据中已有的解答模式，而在分布外（out-of-distribution）变形问题上容易出现性能下降。这一现象在数学竞赛题和算法题等强逻辑任务中尤为明显，严重限制了模型的泛化能力。

\subsection{强化学习后训练的兴起}
\label{subsec:rl_training}

在此背景下，强化学习（Reinforcement Learning，RL）被逐步引入大语言模型的后训练阶段，并成为提升模型复杂推理能力的重要技术路径。相比于SFT的被动学习，RL提供了一种\textit{主动探索}和\textit{奖励优化}的范式。

\textbf{强化学习后训练与SFT的本质区别。} 强化学习后训练（Reinforcement Learning Fine-Tuning，RLFT）的基本思路与监督微调有着本质的不同。SFT要求模型\textit{"模仿答案"，}而RL则要求模型通过与环境的交互来\textit{"探索策略"}。在LLM的语境下，这个\textit{"环境"可以是人类评估者给出的偏好反馈，}也可以是一个自动化的验证器（Verifier）。对于数学、代码等具有客观正确答案的任务，验证器可以由规则系统直接构成：答案是否正确，程序是否通过测试，逻辑链条是否自洽。这类基于可验证奖励的强化学习方法，被称为RLVR（Reinforcement Learning with Verifiable Rewards）\cite{deepseek2024reasoning}，近年来获得了学术界和工业界的广泛关注。

\textbf{策略梯度框架的数学基础。} 在强化学习的框架中，LLM被视为一个策略（Policy），其参数$\theta$决定了在给定输入Prompt $x$下输出回答$a$的概率分布$\pi_\theta(a|x)$。强化学习的目标是最大化策略在所有输入上的期望奖励：

\begin{equation}
J(\theta) = \mathbb{E}_{x \sim d_0, a \sim \pi_\theta(\cdot|x)} [r(x, a)]
\label{eq:rl_objective}
\end{equation}

其中$d_0$是Prompt的采样分布，$r(x, a)$是验证器给出的奖励信号（在可验证任务中通常为二元信号：$r \in \{0, 1\}$）。

通过策略梯度（Policy Gradient）方法\cite{williams1992simple}，我们可以在不知道奖励函数梯度的情况下，利用采样得到的回答对模型参数进行迭代优化：

\begin{equation}
\nabla_\theta J(\theta) \approx \mathbb{E}_{x, a} [A(x, a) \cdot \nabla_\theta \log \pi_\theta(a|x)]
\label{eq:pg_gradient}
\end{equation}

其中$A(x, a)$是优势函数（Advantage Function），衡量回答$a$相对于平均水平的优越程度。

\textbf{RLFT相比SFT的优势。} 这一范式的重要性在于，模型不再完全受限于人类专家示范数据，而可以在输出空间中基于奖励信号进行策略搜索。只要验证器能够提供稳定且可判别的奖励，模型便有可能提升高奖励推理路径的生成概率。换言之，模型可以自主探索和发现不在SFT训练数据中出现过的解题方法。因此，强化学习后训练已成为近年来提升前沿模型推理能力的重要组成部分。相关研究\cite{deepseek2024reasoning}已经证明，RLVR在数学推理、代码生成等任务上能够显著超越纯SFT的基线。

\subsection{规模化训练中的计算代价}
\label{subsec:computational_cost}

然而，强化学习后训练通常伴随显著更高的计算成本。与监督微调相比，RLFT的额外开销主要来自在线采样（On-policy Sampling）的必要性。

\textbf{在线采样的必要性与成本。} 在策略梯度方法中，模型必须基于当前策略$\pi_\theta$持续生成新的样本，用于估计梯度方向。这是因为策略梯度定理（Policy Gradient Theorem）\cite{sutton1999policy}要求所使用的样本必须来自于当前策略的分布——若使用旧策略生成的样本（Off-policy），梯度估计将产生严重的偏差，导致训练不稳定甚至发散。这就意味着，每一轮梯度更新之前，都必须进行一轮代价高昂的推理过程（即让模型生成文本），才能获得用于计算梯度的原始数据。

在实际测量中，对于主流的数学推理任务，单轮训练步骤中推理采样所消耗的计算量通常占整体训练算力的80\%以上。这意味着，训练效率的瓶颈并不在于参数的反向传播，而在于如何高效地从当前策略中获取有价值的样本。换言之，采样阶段已经成为整个RLFT流程中的主要算力消耗者。

\textbf{采样规模与方差-偏差权衡。} 更进一步，主流算法如GRPO（Group Relative Policy Optimization）\cite{gao2024grpo}为了估计每个Prompt的相对优势，通常对每个输入生成$n$个候选回答（$n$一般设置为8至32之间），通过组内相对奖励来计算优势函数。这种\textit{"以采样换精度"的策略}虽然能够降低方差、稳定训练，却同时也将推理计算的开销再次放大了$n$倍。

具体而言，如果单次Prompt的推理时间为$t_\text{inference}$，那么RLFT每步的总推理时间约为$B \times n \times t_\text{inference}$，其中$B$是批次大小。这意味着，即使仅仅将采样数$n$从8增加到16，推理时间也会翻倍——这在实际训练中是一个极其昂贵的权衡。

\textbf{问题的核心：采样效率与计算预算的矛盾。} 因此，如何在保证训练信号质量的前提下，提高单位计算资源对应的有效梯度贡献，已成为LLM强化学习后训练中的关键问题。这也是本文研究的直接动机所在。

\section{核心问题：信号坍塌与采样低效}
\label{sec:signal_collapse}

本节深入分析GRPO框架下的一个关键现象：信号坍塌（Signal Collapse），这是限制采样效率的根本瓶颈。

\subsection{GRPO框架与信号坍塌现象}
\label{subsec:grpo_signal_collapse}

在GRPO框架下，对于一个给定的Prompt $x$，模型生成$n$个回答$\{a_1, a_2, \ldots, a_n\}$，验证器给出对应的奖励向量$\mathbf{r} = [r_1, r_2, \ldots, r_n]$。每个回答的优势函数被定义为：

\begin{equation}
A_\text{GRPO}(x, a_i) = \frac{r_i - \bar{r}}{\sigma_r + \epsilon}
\label{eq:grpo_advantage}
\end{equation}

其中$\bar{r} = \frac{1}{n}\sum_i r_i$是组内奖励均值，$\sigma_r = \sqrt{\frac{1}{n}\sum_i(r_i - \bar{r})^2}$是组内奖励标准差，$\epsilon$是一个防止除零的小量（通常设为$10^{-8}$）。

由上式可知，当$n$个回答全部正确（$r_i = 1$ for all $i$）或全部错误（$r_i = 0$ for all $i$）时，$\sigma_r = 0$，从而$A_\text{GRPO} = 0$，梯度为零，模型参数不发生更新。本文将该现象称为\textit{信号坍塌}（Signal Collapse）。

\textbf{信号坍塌的普遍性。} 信号坍塌并非罕见的边界情形，而是固定采样预算下普遍存在的结构性问题。这一现象可以从两个方面理解：

\begin{enumerate}
    \item \textbf{正向坍塌：} 对于较易题目，随着训练推进，模型生成正确回答的概率逐渐接近1，此时固定$n$次采样更容易全部正确（概率接近$p^n \approx 1$），从而产生$\sigma_r = 0$的情况。
    \item \textbf{负向坍塌：} 对于较难题目，当单次采样命中正确答案的概率$p$很低时，有限预算下更容易出现全部错误的情况（概率约为$(1-p)^n$），同样导致$\sigma_r = 0$。
\end{enumerate}

这两类情形都会使对应Prompt在当前训练步中无法提供有效更新信号，严重浪费了已经消耗的采样计算量。

\textbf{实验证据。} 实验数据表明，该问题在训练过程中具有较高发生比例。在以Qwen2.5-Math-1.5B为基础模型、使用GRPO在数学推理数据集上训练的实验中，统计了不同采样规模$n$下，训练前期与后期中发生信号坍塌的Prompt比例，如表\ref{tab:signal_collapse_rate}所示。

\begin{table}[htbp]
    \centering
    \caption{不同采样规模下的信号坍塌比例}\label{tab:signal_collapse_rate}
    \begin{tabular}{|c|c|c|}
    \hline
    采样规模 $n$ & 训练初期坍塌率 & 训练后期坍塌率 \\
    \hline
    4 & 42.1\% & 78.5\% \\
    8 & 28.3\% & 61.2\% \\
    16 & 15.6\% & 54.8\% \\
    32 & 8.2\% & 50.4\% \\
    \hline
    \end{tabular}
\end{table}

表\ref{tab:signal_collapse_rate}表明，即便将采样规模提升至$n=32$，训练后期仍有超过一半的Prompt无法提供有效更新信号。与此同时，将$n$从4提升到32会使推理计算开销增加8倍，而坍塌率的下降幅度并不与成本增长同比例。这说明，\textbf{仅通过增加固定采样数量缓解信号坍塌，其计算效率是有限的}，这是一个根本性的瓶颈而非可通过简单扩大预算解决的问题。

\subsection{Pass@k实验：上限被遮蔽，而非不存在}
\label{subsec:pass_k_analysis}

信号坍塌现象进一步引出一个深层次的问题：当某个Prompt在训练中频繁发生负向坍塌时，这究竟意味着模型不具备解题能力，还是仅仅意味着有限采样预算不足以显式观测到低概率正确路径？

Pass@$k$实验为这一问题提供了有力的经验依据。在基于Qwen2.5-Math-1.5B的一组系统性实验中，固定模型权重（使用训练中期的checkpoint），分别统计不同采样次数$k$下模型在数学测试集上至少生成一个正确回答的题目比例。结果表明：

\begin{itemize}
    \item 当$k = 1$时，Pass@1仅为26.5\%，意味着大多数题目在单次采样下模型无法给出正确答案。
    \item 当$k = 256$时，Pass@256达到了81.3\%，说明绝大多数题目在充分采样的条件下，模型是具备解题能力的。
\end{itemize}

\textbf{深层含义。} 上述差异（81.3\% vs 26.5\%）说明，在不少情形下，问题并不在于模型完全缺乏解题能力，而在于正确路径虽然已存在于策略分布之中，却因概率质量较低而难以在小样本采样中被观测到。换言之，模型的\textit{能力上界}（由Pass@256反映）与其在\textit{小采样预算下的表现}（由Pass@1反映）之间存在巨大鸿沟，而这个鸿沟正是采样效率不足造成的。

Pass@$k$的数学形式为：
\begin{equation}
\text{Pass}@k = 1 - \left(1 - p\right)^k
\label{eq:pass_at_k}
\end{equation}

其中$p$是模型对该题的单次通过率（内生正确率）。这一指数形式表明，当$p$很小时，需要极大的$k$才能获得显著的Pass@$k$提升。

此外，进一步比较了不同训练阶段checkpoint在$k=4$与$k=256$下的表现差异。在训练较早的checkpoint上，$n=4$时有35.3\%的Prompt触发正向坍塌（全部正确），而在$n=256$时该比例降至10.2\%。这表明，一部分在小采样预算下表现稳定的题目，在更大采样规模下仍可能存在尚未充分覆盖的困难情形。

\textbf{相关文献支持。} 这些观察与已有pass@$k$分析结论具有一致性。例如，Yang等人在近期的研究中指出，RLVR的主要作用更接近于在已有分布内放大高奖励路径的概率质量，而非在基座模型能力边界之外生成全新的推理路径；因此，模型在pass@1上的提升未必对应充分采样条件下能力上界的同步扩展。

本文认为，上述现象的关键约束之一在于主流训练范式中的采样效率不足。在标准GRPO训练中，固定且有限的采样预算使大量低概率但真实存在的正确路径长期处于统计上难以观测的状态，模型因而无法稳定获得与这些路径相关的梯度信号，也难以持续提升其概率质量。这一观察构成了本文方法设计的直接动机。

\subsection{负向坍塌概率与有效更新频率分析}
\label{subsec:collapse_probability_analysis}

为进一步量化固定采样预算对难题学习效率的影响，下面对负向坍塌概率及其对应的有效更新频率进行详细的数学分析。

\textbf{基本概率模型。} 设某Prompt $x$在当前策略下的内生正确率为$p$，即从$\pi_\theta(\cdot|x)$中独立采样的每个回答有概率$p$为正确答案（这里假设采样之间相互独立，这在实践中是合理的近似）。假设我们每轮对该Prompt采样$n$次，则$n$次采样全部错误（负向坍塌）的概率为：

\begin{equation}
P_\text{collapse} = P(K = 0) = (1 - p)^n
\label{eq:collapse_probability}
\end{equation}

其中$K \sim \text{Binomial}(n, p)$是$n$次采样中正确回答的个数。

\textbf{有效训练间隔的定义与计算。} 定义有效训练间隔$T$为平均需要多少轮训练，该Prompt才能获得一次非零梯度更新机会。由于坍塌发生的概率为$P_\text{collapse}$，不坍塌（从而能提供有效梯度）的概率为$1 - P_\text{collapse}$，因此有效训练间隔为：

\begin{equation}
\mathbb{E}[T] = \frac{1}{1 - P_\text{collapse}} = \frac{1}{1 - (1-p)^n}
\label{eq:effective_interval}
\end{equation}

当$p \cdot n \ll 1$时，利用一阶泰勒展开$(1-p)^n \approx 1 - np$，可得近似：

\begin{equation}
\mathbb{E}[T] \approx \frac{1}{np}
\label{eq:effective_interval_approx}
\end{equation}

这个近似在$p$很小的场景下特别有用，它清晰地表明有效训练间隔与$n$成反比、与$p$成反比。

\textbf{具体数值例子。} 当$n$远小于$1/p$时，$\mathbb{E}[T]$将随$n$的不足而快速增大。以一道内生正确率$p = 1/256 \approx 0.0039$的题目为例，若采样数$n = 8$，则：

\begin{align}
P_\text{collapse} &= \left(1 - \frac{1}{256}\right)^8 \approx 1 - \frac{8}{256} = 0.96875 \\
\mathbb{E}[T] &\approx \frac{1}{8 \times 1/256} = 32
\end{align}

这意味着，在总训练步数为600步的训练计划中，该题平均仅能获得$600 / 32 \approx 19$次非零梯度更新机会。若奖励信号本身还存在较强随机性，这样的有效更新次数通常不足以支撑稳定的参数修正。换言之，简单增加训练步数也难以弥补采样效率不足的问题。

\textbf{预算扩张的边界效应。} 进一步地，增大固定采样数$n$对高难题的改善受到真实正确率$p$的显著限制。将$n$从8提升到16时，若$p = 1/256$，则$\mathbb{E}[T]$可由32降至16，这是一个50\%的改善。但若$p = 1/1024$，则在$n=16$时仍有$\mathbb{E}[T] \approx 64$，几乎没有改善。换言之，当$p$足够小时，仅依赖固定预算扩张并不能有效缩短难题的学习间隔，而对应计算成本却近似线性增长。这表明**简单的"暴力采样"（brute-force large-budget sampling）策略在渐进式学习中效率极低**。

\textbf{结论。} 上述分析表明，缓解难题上的信号坍塌问题，关键不在于对所有样本等比例放大$n$，而在于依据样本难度动态分配采样资源。简单粗暴地增加$n$往往会导致简单题上的大量浪费（过采样），而对于极难题目的帮助也有限。

\section{现有解决方案及其不足}
\label{sec:limitations_of_existing_methods}

面对信号坍塌问题，现有研究提出了若干解决方案，但每种方案都存在不可忽视的局限性。本节系统回顾现有方法及其不足，为后续的方法创新奠定基础。

\subsection{被动过滤与固定大规模采样}
\label{subsec:filtering_and_large_budget}

\textbf{被动过滤（Passive Filtering）} 是较为直接的应对策略：在每轮训练中，若某个Prompt的采样结果发生坍塌（即$\sigma_r = 0$），则直接将其从当前batch中剔除，仅利用产生有效信号的Prompt计算梯度。该方法虽然避免了零梯度样本直接进入更新，但会引入训练分布偏移（distribution shift）。随着模型能力提升，简单题更容易因正向坍塌被剔除，难题则更容易因负向坍塌被剔除，最终保留下来的往往是中等难度样本。这可能导致模型在极端困难问题上的学习不足。此外，被动过滤并不减少已发生的推理成本，浪费的采样计算已经被消耗。

\textbf{固定大规模采样（Large Fixed-Budget Sampling）} 是另一类常见思路，即通过将采样数$n$增大至64、128乃至512来降低坍塌概率。该方法在理论上可以逼近充分采样的极限，但计算代价极高。在8张H100 GPU的训练配置下，将$n$从8提升至128会使单步训练时间由约3分钟增至接近50分钟，几乎增加了16倍。这对于实际项目而言是不可承受的，因此难以作为常规训练方案。而且，这种方法无法充分利用模型的能力差异——简单题也被强制采样128次，造成严重浪费。

\subsection{两阶段预算分配与现有自适应采样}
\label{subsec:two_stage_and_adaptive}

\textbf{两阶段预算分配（Two-stage Budget Allocation）} 是一种更精细化的方案：先使用少量样本估计每个Prompt的通过率或难度，再据此分配后续采样预算，将更多资源分配给信息价值最高的题目。然而，这类方法同样存在局限：估计阶段本身需要额外计算资源，小样本条件下的通过率估计也不稳定（容易受到样本波动的影响）；此外，两阶段串行流程与现代推理引擎的流水线并行架构并不完全兼容，工程实现复杂度较高。

\textbf{已有自适应采样工作} 提出，不再为所有Prompt预先设定统一且固定的采样数$n$，而是分轮追加采样，并在满足预设退出条件时提前停止。以代表性工作Reinforce-Ada为例，该类方法通常在每轮采样固定数量样本（如$M=4$）后，判断是否满足有效信号条件，例如\textit{Balance条件}，即至少同时观测到一个正样本和一个负样本。此类机制能够减少简单题上的冗余采样，并为难题保留进一步追加预算的可能。然而，现有自适应采样方法在以下\textbf{三个方面}仍存在进一步改进空间，这也是本文关注的核心问题：

\begin{enumerate}
    \item \textbf{Balance条件在训练后期的额外开销：} 当模型对某些题目的通过率已接近1时，为满足同时观测正负样本的停止条件，往往需要继续采样以寻找少量负样本，从而增加不必要的计算成本。这对应表\ref{tab:signal_collapse_rate}中观察到的"训练后期坍塌率仍然很高"的现象。
    
    \item \textbf{固定小批次与KV Cache并行架构的适配问题：} 若每轮仅采样少量样本（如4个）并立即判断是否退出，则需要频繁重新调度Prefill计算，前缀复用率受限，推理吞吐量难以充分释放。现代推理引擎（如vLLM）的设计目标就是通过KV Cache共享和Prefill共享来提升吞吐，但固定小批次的频繁调度无法充分利用这些优化。
    
    \item \textbf{动态温度与分布偏差的校正问题：} 当采样轮次增加时，适度提高采样温度有助于提升对低概率正确路径的探索能力（这与深度强化学习中的exploration-exploitation权衡密切相关），但温度变化会引入采样分布与策略分布之间的偏差。若不进行校正，梯度估计可能产生系统性偏差，导致收敛性问题。
\end{enumerate}

\textbf{总结：} 总体而言，现有方案尚未同时兼顾采样效率、硬件友好性与梯度估计正确性。本文旨在围绕这三个维度提出一套统一的改进框架。

\section{研究内容与主要贡献}
\label{sec:contributions}

针对上述问题，本文开展了以下\textbf{三个层面}的研究工作：

\subsection{理论层面：难度感知的信息效率分析与退出条件优化}
\label{subsec:theory_contribution}

本文从信息论视角建立了题目难度与单位采样成本信息效率之间的定量关系，证明了难度感知的预算分配相较于难度无关的均匀分配具有更高的信息增益效率。基于此，本文提出以"正信号优先（Positive-first）"退出条件替代"平衡条件（Balance）"作为训练后期的采样停止准则，并给出相应的理论分析与实验验证。

具体而言，本文定义了单位采样成本信息效率函数$\eta(p) = p(1-p)^2$，证明其在$p=1/3$处达到最大值，意味着中等难度的题目（$p \approx 1/3$）在单位采样成本下能提供最高的信息增益。这一结论为采样预算的重新分配提供了理论依据。

\subsection{系统层面：指数增长批次采样策略}
\label{subsec:system_contribution}

本文提出了指数增长批次采样策略（Exponential Batch Sampling，EBS），令累计采样预算按轮次呈几何扩展；在本文默认设置下，对应的逐轮批次为$2, 2, 4, 8, 16$，即初始两轮保持相同规模，此后按2倍增长。该设计在保留自适应退出灵活性的同时，提高了推理引擎中的KV Cache前缀复用率，从而在固定算力预算下缩短单步训练时间。

EBS策略的核心创新在于：
\begin{itemize}
    \item 简单题（$p \approx 1$）在第1-2轮即快速退出，仅消耗少量Prefill和Decode计算。
    \item 困难题（$p \ll 1$）进入多轮采样，后期批次的增大能够有效摊销Prefill成本，使KV Cache命中率接近单轮固定采样。
    \item 这种设计天然实现了难度感知的预算分配，无需显式的难度估计器。
\end{itemize}

\subsection{算法层面：温度分层重要性采样机制}
\label{subsec:algorithm_contribution}

本文提出层级几何自适应采样框架（Hierarchical Geometric Adaptive Sampling，H-GAS）中的温度分层重要性采样（Temperature-Stratified Importance Sampling，TSIS）机制，将动态温度缩放（Dynamic Temperature Scaling）与重要性采样校正统一到同一套策略梯度框架之中。

具体而言，本文令采样温度随累计采样深度对数增长（$T_t = T_0 \cdot (1 + \alpha \log_2(N_t / M_0))$），以提升深度采样阶段对低概率正确路径的探索能力；同时，对来自不同温度的样本统一引入重要性权重，对策略梯度项与基准线估计进行联合校正。

TSIS机制的关键意义在于：
\begin{itemize}
    \item 重要性权重不仅用于修正多温度采样带来的分布偏差，也隐式刻画了样本相对于目标策略的"可信度"与"训练价值"。
    \item 低温下高置信生成但最终错误的样本会获得更强的惩罚信号，这些"有信心的失败"对模型的改进最为直接。
    \item 高温下偶然探索到的正确样本则会被适度折现（通过较低的$\rho_i$），避免训练过程过度偏向侥幸解。
\end{itemize}

\subsection{综合框架}
\label{subsec:framework_integration}

上述三项贡献共同构成层级几何自适应采样框架H-GAS：
\begin{itemize}
    \item Positive-first对应解决第\ref{sec:limitations_of_existing_methods}节中的"Balance条件后期额外开销"问题。
    \item EBS对应解决"固定小批次与KV Cache不适配"问题。
    \item TSIS对应解决"温度变化引入分布偏差未校正"问题。
\end{itemize}

整体框架兼顾采样效率、系统吞吐与统计正确性，并在Qwen2.5-Math-1.5B、Qwen2.5-Math-7B等开源数学推理模型上进行了实验评估。相比于现有方法，H-GAS实现了\textbf{三维优化}：
\begin{enumerate}
    \item \textbf{采样效率维度}：通过Positive-first和EBS，在相同采样预算下能服务更多的有效Prompt，避免无效坍塌浪费。
    \item \textbf{硬件利用维度}：通过指数增长的批次结构，更好地适配现代推理引擎的KV Cache机制，提升实际训练速度。
    \item \textbf{统计正确性维度}：通过TSIS机制，在利用动态温度探索的同时，确保梯度估计的无偏性和收敛性。
\end{enumerate}

\section{论文组织结构}
\label{sec:organization}

本论文共六章，各章内容安排如下：

\begin{description}
    \item[第一章 引言] 介绍研究背景、核心问题、现有方法局限与本文贡献。本章通过系统分析大语言模型强化学习后训练中的采样效率瓶颈，为后续工作奠定基础。
    
    \item[第二章 相关工作综述] 梳理大语言模型强化学习后训练的技术演进、采样效率优化方向的现有研究，以及推理引擎并行加速技术的相关工作。通过全面回顾现有文献，明确本文的创新定位。
    
    \item[第三章 方法论] 系统阐述本文提出的层级几何自适应采样框架，包括难度感知的信息效率分析、指数增长批次策略的设计原理，以及动态温度缩放与重要性采样校正的算法推导。本章提供了完整的理论和算法支撑。
    
    \item[第四章 实验] 介绍实验设置，并分别从主要性能对比、模块消融分析、计算效率验证及训练动态分析四个维度呈现实验结果。通过多角度的实验验证，证明H-GAS的有效性。
    
    \item[第五章 讨论] 深入分析各方法模块背后的作用机制，并指出本文研究的局限性与未来改进方向。这一章为读者提供了更深层的理解和思考空间。
    
    \item[第六章 结论] 总结全文的主要工作和贡献，并展望后续研究方向。
\end{description}

通过这样的组织结构，本文形成了一个完整的论证链条：从问题识别（第一章）、文献回顾（第二章）、方法设计（第三章）、到实验验证（第四章）和深入讨论（第五章），最后总结贡献（第六章）。
